\documentclass{article}
\usepackage{PRIMEarxiv} % Core package for the PrimeAI Template
\usepackage{amsmath, amssymb, amsthm, bm, dcolumn} % Add amsthm here for the proof environment
\usepackage[numbers,sort&compress]{natbib} % Natbib for citations
\usepackage{graphicx} % For high-quality images
\usepackage{multicol} % Multi-column figures/tables
\usepackage{hyperref} % Hyperlinks
\usepackage{listings} % Code listings
\usepackage{authblk} % For structured affiliations
% Define theorem style (optional)
\theoremstyle{plain}
\newtheorem{theorem}{Theorem}
\renewcommand{\qedsymbol}{}

% Custom commands for primes and qubit encoding
\newcommand{\primeQubit}[1]{|p_{#1}\rangle}
\newcommand{\primeGate}[1]{U_{p_{#1}}}

\usepackage{wrapfig}
\usepackage[pscoord]{eso-pic}
\usepackage[fulladjust]{marginnote}
\reversemarginpar

% Typesetting improvements without footnote patching
\usepackage[protrusion=true, expansion=true, tracking=false]{microtype}
\microtypecontext{spacing=nonfrench}

% Line numbers
\usepackage[right]{lineno}

% Text layout - adjust as needed
\raggedright
\setlength{\parindent}{0.5cm}
\textwidth 5.25in 
\textheight 8.75in

% Set double spacing
\usepackage{setspace} 
\doublespacing

% Adjust width for specific content
\usepackage{changepage}

% Adjust caption style
\usepackage[aboveskip=1pt,labelfont=bf,labelsep=period,singlelinecheck=off]{caption}

% Remove brackets from references
\makeatletter
\renewcommand{\@biblabel}[1]{\quad#1.}
\makeatother

% Header, footer, and page numbers
\usepackage{lastpage,fancyhdr}
\pagestyle{fancy}
\fancyhf{}
\fancyhead[L]{Preprint - PrimeAI Enhanced Template}
\fancyfoot[C]{\scriptsize Multiplicity Theory © 2024 Ryan Van Gelder - Citizen Gardens \\ Licensed Under MIT and CC BY-NC-SA 4.0.}
\fancyfoot[R]{Page \thepage\ of \pageref{LastPage}}
\renewcommand{\footrule}{\hrule height 2pt \vspace{2mm}}

\begin{document}
\title{Quantum Folklore: Exploring Archetypes and Myths through Multiplicity}

\author{Ryan Van Gelder}
\affil{Citizen Gardens - The Foundation of Multiplicity}
\date{\today}
\maketitle

\begin{abstract}
This paper explores the application of Multiplicity Theory to the study of cultural archetypes and myths, proposing a framework that models these as dynamic, interconnected systems exhibiting quantum-like transitions. By leveraging prime-based encoding, recursive feedback mechanisms, and tensor network representations, we aim to map the evolution and entanglement of myths across cultures, creating what we term a \textit{quantum folklore}. This approach offers novel insights into the interconnectedness and adaptability of human storytelling, emphasizing the non-linear dynamics and emergent behaviors inherent in cultural narratives.
\end{abstract}

\section{Introduction}
Cultural archetypes and myths have long served as repositories of collective human experience, transcending temporal and geographical boundaries. Traditional studies often treat myths as isolated constructs; however, Multiplicity Theory  invites a holistic perspective, where myths are interconnected elements within a multi-dimensional framework. This paper introduces a computational model inspired by quantum mechanics to map and analyze these connections.

\section{Multiplicity Framework for Myths}

\subsection{Prime-Based Encoding of Archetypes}
Each archetype $A_i$ is assigned a unique prime identifier $p_i$:\newline
\begin{equation}
\phi(A_i) = p_i,
\end{equation}
where $\phi$ is the encoding function. Myths, as combinations of archetypes, are represented by their prime product:
\begin{equation}
M = \prod_{i=1}^N p_i^{\alpha_i},
\end{equation}
where $\alpha_i$ indicates the frequency or significance of archetype $A_i$ within the myth $M$. This encoding preserves the uniqueness of myths while enabling efficient computational analysis.

\subsection{Dynamic Multiplicity Equation}
The evolution of archetypes across cultures is governed by the Dynamic Multiplicity Equation \cite{galioto2024dynamic}:
\begin{equation}
\frac{\partial \rho_k}{\partial t} = \alpha_k(t)\rho_k + \beta_k(t)I_k + \gamma_k(t) \sum_{j} T_{kj}\rho_j + \lambda(t)\left(\Omega_B(\rho) + \Omega_{FS}(\rho)\right),
\end{equation}
where:
\begin{itemize}
    \item $\rho_k$: The state of archetype $A_k$ at time $t$.
    \item $T_{kj}$: Tensor coefficients capturing interactions between archetypes.
    \item $\Omega_B$ and $\Omega_{FS}$: Geometric terms representing feedback loops from cultural and social systems.
\end{itemize}

\section{Quantum-Like Transitions in Myths}

\subsection{Entanglement of Archetypes}
The entanglement of two archetypes $A_i$ and $A_j$ is modeled as a correlation in their quantum states:
\begin{equation}
\psi(t) = \sum_{i,j} C_{ij}(t) \phi(A_i) \phi(A_j),
\end{equation}
where $C_{ij}(t)$ represents the time-dependent entanglement coefficient. This formalism captures shared themes or symbolic overlaps between myths.

\subsection{Superposition of Mythological States}
Archetypes can exist in superposition, embodying multiple cultural interpretations simultaneously:
\begin{equation}
\Phi = \sum_{i} \alpha_i \phi(A_i),
\end{equation}
where $\alpha_i$ are the amplitudes corresponding to each archetype's contribution.

\section{Tensor Networks for Interconnected Narratives}
To model the multi-dimensional dependencies between myths, we employ tensor networks:\newline
\begin{equation}
T = \sum_{i,j,k} T_{ijk} \phi(A_i) \phi(A_j) \phi(A_k),
\end{equation}
where $T_{ijk}$ encodes the strength of the interaction between archetypes $A_i$, $A_j$, and $A_k$. These networks facilitate the visualization and analysis of complex cultural narratives.

\section{Applications and Future Directions}

\subsection{Predicting Evolution of Myths}
Using recursive feedback loops, the framework predicts potential evolutions of myths by modeling cultural reinterpretations:\newline
\begin{equation}
M(t+1) = f(M(t), \rho(t)),
\end{equation}
where $f$ incorporates historical and social dynamics.

\subsection{Cross-Cultural Symmetries}
By examining eigenvalues of adjacency matrices derived from tensor networks, we identify universal archetypes and their emergent properties.

\subsection{Educational Impact}
The quantum folklore model fosters appreciation for cultural interconnectedness, aiding educational and interdisciplinary initiatives.

\section{Conclusion}
Multiplicity Theory provides a transformative lens for understanding myths and archetypes as dynamic systems. By integrating quantum-inspired principles, this framework elucidates the non-linear and emergent behaviors of cultural narratives, paving the way for deeper interdisciplinary research.

\nocite{*}
\bibliographystyle{plain}
\bibliography{references}

\end{document}